\documentclass[12pt]{article}
\usepackage[a5paper,margin=2cm]{geometry}
\usepackage{lmodern}
\usepackage[T1]{fontenc}
\usepackage{amsmath}
\usepackage[style=numeric,backend=biber]{biblatex}
\usepackage{hyperref}
\addbibresource{refs.bib}

\title{Lantern build times: appendix}
\author{A. Learner}
\date{}

\begin{document}
\maketitle
\appendix

\section{Every run}
\label{sec:runs}

The paper reports the median of five builds for each folder. The table
below lists all fifteen builds, in seconds, in the order they ran.

\begin{table}[h]
  \centering
  \begin{tabular}{rrrrrr}
    \hline
    Notes & Run 1 & Run 2 & Run 3 & Run 4 & Run 5 \\
    \hline
    100 & 0.21 & 0.20 & 0.20 & 0.19 & 0.20 \\
    1000 & 1.93 & 1.88 & 1.90 & 1.91 & 1.87 \\
    10000 & 19.6 & 19.4 & 19.2 & 19.4 & 19.5 \\
    \hline
  \end{tabular}
  \caption{Build time of every run, in seconds.}
  \label{tab:runs}
\end{table}

\section{Setup}
\label{sec:setup}

All builds ran on one laptop with four cores and sixteen gigabytes of
memory, plugged in, with the network disconnected so that no
background update could start in the middle of a run.

Before the first build of each folder, every note was read once, so
that the files were in the operating system's cache for all five runs.
Without that, the first run of the largest folder took almost twice as
long as the others, which says more about the disk than about Lantern.

The mesurements were taken with the shell's own timer, which reports
the wall-clock time of the whole command, from start-up to the last
page written.

\printbibliography
\end{document}
